\scanpage{41}
\mspage{39}
% archive sheet number: 40
peut supposer $\unsure{U'_\alpha|V_\alpha}=U_\alpha$ (EGA \uline{IV} 8), donc $\unsure{E_\alpha|\mathscr E_\alpha}$ est un faisceau constant tordu sur $U_\alpha$, \inserted{donc provient d'une représentation $\ell$-adique de $\pi_\alpha=\pi_1(U_\alpha,\overline\eta)$ dans $\unsure{\operatorname{Aut}E}$.} En vertu du th. de changement de base propre (SGA 4 \uline{XII}), $Rf_{K*}((\unsure{\mathbb Z_\ell|Z_\ell})_{X_K})=E_K$ est \deleted{\unsure{lui}} isomorphe à l'image inverse de $E_{U_\alpha}$ par \deleted{\unsure{de $g$}} $\eta\longrightarrow U_\alpha$, donc la représentation de $\pi$ \deleted{\unsure{provenant}} définie par $H^i(X_{\overline K},\mathbb Z_\ell)$ est celle déduite de $\pi_\alpha$ par l'hom canonique $\pi-\pi_\alpha$.
\margin{\unsure{donc est} $E$ de t.f. sur $\mathbb Z_\ell$}
\[
\deleted{\pi\longrightarrow\pi_\alpha}.
\]
\deleted{Notons que $\pi_\alpha$ est une extension de $\pi_{\alpha,0}=\operatorname{Gal}(\overline k_\alpha/k_\alpha)$}

Pour la suite, mettons l'exposant $'$ au lieu de l'indice $\alpha$, et écrivons
\[
\unsure{\varphi|\alpha}\qquad \pi-\pi'.
\]
Notons que $\pi'$ est une extension de
\[
\pi'_0=\operatorname{Gal}(\overline{k'}/k')\simeq\pi_1(S'_{\deleted{\unsure{\alpha}}})
\]
par un ss-groupe \deleted{\unsure{de}} $I'$, et que l'hom $\varphi$ induit un hom :
\[
I\longrightarrow I'\subset\pi'.
\]
Dans le cas où $V_\alpha$ est de val. discrète, \deleted{\unsure{par}} la conclusion résulte de n° 6.
% slanted
\margin{Ceci \unsure{suffit|suffirait} déjà à entraîner le cas 8.2 a), qui ne sera pas \unsure{démontré} pour l'instant par la dém. ``géom.'' de l'exposé suivant (\unsure{ce nom est|ceci étant} \unsure{la seule application|la seule appellation indiquée} pour le moment de la méthode ``arithmétique'', \unsure{le lecteur pourra|la lecture permettant} \unsure{se dispenser de lire|de dispenser le lecteur, si} \unsure{le cas} \uuline{général} \unsure{le gêne} !)}

\deleted{Nous allons utiliser donc le théorème de structure suivant, qui est une \inserted{conséquence} \unsure{variante} du ``lemme de Abhyankar'' :}

\deleted{\uline{Lemme} 8.5. Soit $A$ un anneau local noethérien \uline{hensélien}, $k$ le corps résiduel \unsure{$k$}, $p$ l'exp. caractéristique de $k$, \unsure{Spec} $S=\operatorname{Spec}A$, $D=\sum_{i\in I}D_i$ un diviseur à croisements normaux, $(x_i)_{i\in I}$ une partie d'un syst. de paramètres, $a=\prod x_i$, $U=S_a$, $\unsure{\xi|\zeta}$ un pt géom. de $U$. Alors}
\deleted{Alors $\xi$ définit une \deleted{extension séparable} \inserted{clôture séparable} $\overline k$ de $k$, et on}
